% ============================================================
% M3 S3 练习件·W1 臂B 片件 —— 课时04 点斜式与斜截式（2.2.2上）＝母版照抄件
% 生成：M3 成卷轮 S3 W1 臂B｜2026-09-14｜编译 cwd＝本目录（中文路径禁 \input 绝对路径）
%   xelatex main-true.tex ×2 → 含详解印本；xelatex main-false.tex ×2 → 纯题版（单源双本）
% 输入链（全只读）：题面库 课时04-点斜式与斜截式.md（题面逐字装配；【注】行＝装配层注记不印，
%   题首「（预习注：…）」行＝印面照录）＋ -答案侧.md（值逐字回嵌）
%   ＋ 定稿/批A-课时04-点斜式与斜截式.md（详解回嵌源）；toolchain＝qp-m3.sty（钉后版
%   md5 7c3930362be8a0a2bdf21bbf8ac16573，本地挂载副本，破损即停）。
% 母版体例：照抄 成卷/练习件/课时01-坐标法/《练习件母版体例说明.md》；骨架与门谱原样，只换内容层。
% 键账：件内 ansblock 键集 ≡ 题面库片练键集 ≡ manifest 练键序 E1~E16（16 键，零缺漏零多余）；
%   拓区隔离：课时04 拓槽 T1 键不入本件（S4 拓展册收）；导学键 G1~G5 属导学件域，亦不入。
% 卷式例外案B：练习件不属卷式——答案块物理落点恒为题后 ansblock，禁卷末附卷制；
%   全线括线模（导言 \ansblockgrayfalse 一次置定，件内不切换；灰底盒不可跨栏断，练习件禁用）。
% 题序＝manifest 键序 E1~E16（零重排）；印面题号 1~16 连号（\tihao 号＝\ansitem 号同键）；
%   三档切位 1~7／8~14／15~16（照库序切位，非难度重排；E11~E16＝0913收口§六自课时05前移迁入席，
%   预习注随迁循 04-E4 先例、详解括注随迁四处均照题面库冻面照录）。
% 槽配实录：单选6（E1/E2/E3/E11/E12/E13）＋填空2（E4/E14）＋解答8（E5~E10/E15/E16）；
%   10选4填2解 为波次方案基线槽配，本片实配照题面侧头标（批A §四 满足度表同口径；多选0）。
% 选项槽位：默认四连排（0.25\linewidth−0.25\lxhang），超宽降档梯 四连排→两连排
%   （0.5\linewidth−0.5\lxhang−0.25em）→单列 \lxopt；禁缩字号/负 kern/删标点凑宽。
%   本片实测降档：题2 单列、题3/题13 两连排，余三道（题1/题11/题12）四连排落位（登记见值台账＋回执）。
% 解答书写区：E5/E6 无小问 \liubai[16mm]；E7/E8/E16 两问 \liubai[24mm]；E9/E10/E15 三问
%   \liubai[32mm]（默认16mm·两问24mm·三问以上32mm·cap 40mm；纯题档吞 ansblock 不吞 \liubai）。
% 填空印答：E4/E14 空位 \kongda{值}（详解版印答、纯题版退定宽空线，两档等形）。
% ============================================================
\documentclass[fontset=none]{ctexart}
\usepackage{qp-m3}
% —— 印面逐字符号路由补钉（承导学件母版；− 走 TNR、▱ NSC 子块；禁改 sty 保 md5） ——
\xeCJKDeclareCharClass{Default}{"2212}%
\setCJKfamilyfont{nscblk}[Path=C:/提示词/工作区/字替对照-0909/variantF/fonts/,BoldFont=NSC-w600.ttf]{NSC-w500.ttf}%
\xeCJKDeclareSubCJKBlock{nscblk}{"25B1}%
\newcommand{\pxparallelogram}{{\CJKfamily{nscblk}▱}}%
% —— 断行微调（承导学件母版实测档，三零门承重；\emergencystretch 不另设） ——
\xeCJKsetup{CJKglue={\hskip 0.10em plus 0.08em minus 0.05em}}%
% —— 练习件全线括线模（S3 波次方案 §四.3：导言一次置定、件内不切换） ——
\ansblockgrayfalse
% —— 页脚件名（qp-headfoot 复用入口） ——
\renewcommand{\qpzhangming}{第二章\quad 平面解析几何}%
\renewcommand{\qpjianming}{练习件}%

\begin{document}
%装配轮S6三本跨本连续页码（G7方案B）setcounter 注入位
\setcounter{page}{84}

\keshi{第4课时\quad 点斜式与斜截式}
\begin{multicols}{2}
\emergencystretch=1em

% ============================================================
% 基础巩固（题1~7＝E1~E7：单选3＋填空1＋解答3；题后紧跟答案制，全线括线 ansblock）
% ============================================================
\dangbadge{基}{础}{巩}{固}

\tihao{1}\zu{直线方程形式与互化×单选} \tieside{简单(知识点二)}直线\(y-b=2(x-a)\)在\(y\)轴上的截距为\nobreak(\kongwei)
\bindopt {\fontsize{10.5pt}{\qpTJgLeadLiti}\selectfont \makebox[\dimexpr0.25\linewidth-0.25\lxhang\relax][l]{A．\(a+b\)}\makebox[\dimexpr0.25\linewidth-0.25\lxhang\relax][l]{B．\(2a-b\)}\makebox[\dimexpr0.25\linewidth-0.25\lxhang\relax][l]{C．\(b-2a\)}\makebox[\dimexpr0.25\linewidth-0.25\lxhang\relax][l]{D．\(|2a-b|\)}\par}

\begin{ansblock}[2章-练-课时04-E1]
% ans:2章-练-课时04-E1
\ansitem{1}{C}
\ansnote{详解}{将\(y-b=2(x-a)\)化为斜截式：\(y=2x-2a+b\)，故直线在\(y\)轴上的截距为\(b-2a\)（\(x=0\)时的纵坐标值，可正可负，不加绝对值）．选：C．}
\end{ansblock}

\tihao{2}\zu{直线方程形式与互化×单选} \tieside{中档(知识点三)}若直线\(l\)经过点\(P(2,3)\)，且在\(x\)轴上的截距的取值范围是\((-1,3)\)，则其斜率\(k\)的取值范围是\nobreak(\kongwei)

\lxopt{A．\((-\infty,-3)\cup(1,+\infty)\)}

\lxopt{B．\((-1,1/3)\)}

\lxopt{C．\((-3,1)\)}

\lxopt{D．\((-\infty,-1)\cup(1/3,+\infty)\)}

\begin{ansblock}[2章-练-课时04-E2]
% ans:2章-练-课时04-E2
\ansitem{2}{A}
\ansnote{详解}{设\(l\)与\(x\)轴交于\(M_0(t,0)\)，则\(t\in(-1,3)\)且\(t\neq2\)（\(t=2\)时直线过\(P\)与\(x\)轴交点同横坐标，即\(x=2\)，斜率不存在，截距为\(2\in(-1,3)\)——此情形\(k\)不存在，题设“斜率\(k\)”已默认存在，排除）．\(k=(3-0)/(2-t)=3/(2-t)\)．当\(t\in(-1,2)\)时\(2-t\in(0,3)\)，\(k=3/(2-t)\in(1,+\infty)\)；当\(t\in(2,3)\)时\(2-t\in(-1,0)\)，\(k\in(-\infty,-3)\)．故\(k\in(-\infty,-3)\cup(1,+\infty)\)．选：A．（临界图法同件源：取\(x\)轴上两端点\(M(-1,0)\)、\(N(3,0)\)，\(k_{PM}=3/3=1\)，\(k_{PN}=3/(-1)=-3\)，\(l\)与开线段\(MN\)（挖去\((2,0)\)上方对应竖线）相交即\(k>1\)或\(k<-3\)．）}
\end{ansblock}

\tihao{3}\zu{直线方程形式与互化×单选} \tieside{中档(知识点二)}已知直线\(ax+by-1=0\)在\(y\)轴上的截距为\(-1\)，且它的倾斜角是直线\(\sqrt{3}x-y-\sqrt{3}=0\)的倾斜角的2倍，则\(a\)，\(b\)的值分别为\nobreak(\kongwei)
\bindopt {\fontsize{10.5pt}{\qpTJgLeadLiti}\selectfont \makebox[\dimexpr0.5\linewidth-0.5\lxhang-0.25em\relax][l]{A．\(\sqrt{3}\)，1}\makebox[\dimexpr0.5\linewidth-0.5\lxhang-0.25em\relax][l]{B．\(\sqrt{3}\)，\(-1\)}\par}
\bindopt {\fontsize{10.5pt}{\qpTJgLeadLiti}\selectfont \makebox[\dimexpr0.5\linewidth-0.5\lxhang-0.25em\relax][l]{C．\(-\sqrt{3}\)，1}\makebox[\dimexpr0.5\linewidth-0.5\lxhang-0.25em\relax][l]{D．\(-\sqrt{3}\)，\(-1\)}\par}

\begin{ansblock}[2章-练-课时04-E3]
% ans:2章-练-课时04-E3
\ansitem{3}{D}
\ansnote{详解}{直线\(\sqrt{3}x-y-\sqrt{3}=0\)的斜率为\(\sqrt{3}\)，倾斜角\(\alpha\)满足\(\tan\alpha=\sqrt{3}\)，\(\alpha=60^\circ\)；所求直线倾斜角\(\beta=2\alpha=120^\circ\)，斜率\(k=\tan120^\circ=-\sqrt{3}\)；又在\(y\)轴上的截距为\(-1\)，故斜截式方程为\(y=-\sqrt{3}x-1\)，即\(\sqrt{3}x+y+1=0\)，比对\(ax+by-1=0\)的常数项须为\(-1\)，须将\(\sqrt{3}x+y+1=0\)两边乘\(-1\)得\(-\sqrt{3}x-y-1=0\)，所以\(a=-\sqrt{3}\)，\(b=-1\)．选：D．}
\end{ansblock}

\tihao{4}\zu{直线方程形式与互化×填空} \tieside{中档(知识点二)}直线\(l\)在\(x\)轴上的截距比在\(y\)轴上的截距大1，且过定点\(A(6,-2)\)，则直线\(l\)的方程为\kongda{x＋2y−2＝0或2x＋3y−6＝0}．

\begin{ansblock}[2章-练-课时04-E4]
% ans:2章-练-课时04-E4
\ansitem{4}{x＋2y−2＝0或2x＋3y−6＝0}
\ansnote{详解}{设直线在\(x\)轴上的截距为\(a\)（\(a\neq0\)且\(a-1\neq0\)），则在\(y\)轴上的截距为\(a-1\)，由截距式：\(x/a+y/(a-1)=1\)．\par\noindent 将\((6,-2)\)代入：\(6/a-2/(a-1)=1\)，去分母：\(6(a-1)-2a=a(a-1)\)，即\(4a-6=a^{2}-a\)，\(a^{2}-5a+6=0\)，解得\(a=2\)或\(a=3\)．\par\noindent \(a=2\)时：\(x/2+y/1=1\)，即\(x+2y-2=0\)；\(a=3\)时：\(x/3+y/2=1\)，即\(2x+3y-6=0\)．\par\noindent （另须检验截距为0的情形：若过原点，设\(y=mx\)，代\((6,-2)\)得\(m=-1/3\)，此时两截距都为0，差为\(0\neq1\)，不合．）所以直线\(l\)的方程为\(x+2y-2=0\)或\(2x+3y-6=0\)．}
\end{ansblock}

\tihao{5}\zu{直线方程形式与互化×解答} \tieside{简单(知识点一)}判断下列各点是否在直线\(y=-6x+7\)上：\(A(0,7)\)，\(B(2,5)\)，\(C(1,1)\)．
\liubai[16mm]{解答书写区}

\begin{ansblock}[2章-练-课时04-E5]
% ans:2章-练-课时04-E5
\ansitem{5}{A在，B不在，C在．}
\ansnote{详解}{\(x=0\)时\(y=-6\times0+7=7\)，故\(A(0,7)\)在直线上；\par\noindent\(x=2\)时\(y=-12+7=-5\neq5\)，故\(B(2,5)\)不在直线上；\par\noindent\(x=1\)时\(y=-6+7=1\)，故\(C(1,1)\)在直线上．}
\end{ansblock}

\tihao{6}\zu{直线方程形式与互化×解答} \tieside{简单(知识点一)}已知点\((a,1)\)，\((-1,b)\)确定的直线方程是\(y=-3x+1\)，求\(a\)，\(b\)的值．
\liubai[16mm]{解答书写区}

\begin{ansblock}[2章-练-课时04-E6]
% ans:2章-练-课时04-E6
\ansitem{6}{a＝0；b＝4}
\ansnote{详解}{两点都满足直线方程：\(1=-3a+1\)，得\(a=0\)；\(b=-3\times(-1)+1=4\)．\par\noindent（检验：\(a=0\)时两点为\((0,1)\)、\((-1,4)\)，确为\(y=-3x+1\)上两点，两点横坐标不等，直线唯一．）}
\end{ansblock}

\tihao{7}\zu{直线方程形式与互化×解答} \tieside{简单(知识点一)}根据下列条件求直线的点斜式方程：(1) 经过点\(A(2,5)\)，斜率为4；(2) 经过点\(B(-2,2)\)，倾斜角为\(30^\circ\)．
\liubai[24mm]{解答书写区}

\begin{ansblock}[2章-练-课时04-E7]
% ans:2章-练-课时04-E7
\ansitem{7}{(1) y−5＝4(x−2)；(2) y−2＝(√3/3)(x＋2)}
\ansnote{详解}{(1) 由点斜式\(y-y_0=k(x-x_0)\)，得\(y-5=4(x-2)\)．\par\noindent (2) 斜率\(k=\tan30^\circ=\sqrt{3}/3\)，故\(y-2=(\sqrt{3}/3)[x-(-2)]=(\sqrt{3}/3)(x+2)\)．}
\end{ansblock}

% ============================================================
% 综合提升（题8~14＝E8~E14：解答3＋单选3＋填空1，照库序切位；题后紧跟答案制）
% ============================================================
\dangbadge{综}{合}{提}{升}

\tihao{8}\zu{直线方程形式与互化×解答} \tieside{简单(知识点二)}根据下列条件求直线的斜截式方程：(1) 斜率是\(\sqrt{3}/2\)，截距是\(-2\)； (2) 倾斜角是\(135^\circ\)，截距是3．
\liubai[24mm]{解答书写区}

\begin{ansblock}[2章-练-课时04-E8]
% ans:2章-练-课时04-E8
\ansitem{8}{(1) y＝(√3/2)x−2；(2) y＝−x＋3}
\ansnote{详解}{(1) \(k=\sqrt{3}/2\)，\(b=-2\)，由斜截式\(y=kx+b\)得\(y=(\sqrt{3}/2)x-2\)．\par\noindent (2) \(k=\tan135^\circ=-1\)，\(b=3\)，故\(y=-x+3\)．}
\end{ansblock}

\tihao{9}\zu{直线方程形式与互化×解答} \tieside{中档(知识点二)}根据下列直线方程，分别求出直线的斜率及在\(y\)轴上的截距：(1) \(y-2=x+1\)； (2) \(y+4=\sqrt{3}(x-2)\)； (3) \(4x+y-3=0\)．
\liubai[32mm]{解答书写区}

\begin{ansblock}[2章-练-课时04-E9]
% ans:2章-练-课时04-E9
\ansitem{9}{(1) k＝1，b＝3；(2) k＝√3，b＝−2√3−4；(3) k＝−4，b＝3}
\ansnote{详解}{(1) \(y-2=x+1\)化为\(y=x+3\)，故\(k=1\)，截距\(b=3\)（定点\((-1,2)\)）；\par\noindent (2) 化为\(y=\sqrt{3}x-2\sqrt{3}-4\)，故\(k=\sqrt{3}\)，截距\(b=-2\sqrt{3}-4\)；\par\noindent (3) 化为\(y=-4x+3\)，故\(k=-4\)，截距\(b=3\)．}
\end{ansblock}

\tihao{10}\zu{直线方程形式与互化×解答} \tieside{中档(知识点二)}判断下列命题的真假：(1) 如果直线\(l\)的斜率不存在，则直线\(l\)在\(x\)轴上的截距唯一；(2) 如果直线\(l\)的斜率存在，则直线\(l\)的截距唯一；(3) 存在直线\(l\)，使得\(l\)在\(x\)轴上的截距与在\(y\)轴上的截距都为0．
\liubai[32mm]{解答书写区}

\begin{ansblock}[2章-练-课时04-E10]
% ans:2章-练-课时04-E10
\ansitem{10}{(1) 真；(2) 假；(3) 真}
\ansnote{详解}{(1) 斜率不存在时\(l\)为竖直线\(x=m\)，它与\(x\)轴恰有一个交点\((m,0)\)，\(x\)轴截距为\(m\)，唯一，真命题；\par\noindent (2) 斜率存在时\(l\)与\(y\)轴必只有一个交点（\(y\)轴截距唯一），但与\(x\)轴未必相交：水平直线\(y=b\)（\(b\neq0\)，如\(y=1\)）与\(x\)轴无交点，\(x\)轴截距不存在，故“截距唯一”为假命题；\par\noindent (3) 过原点的直线（如\(y=x\)）与两轴交点都为原点\((0,0)\)，两截距都为0，真命题．}
\end{ansblock}

\tihao{11}\zu{直线方程形式与互化×单选} \tieside{简单(知识点三)}经过\(M(3,2)\)与\(N(6,2)\)两点的直线的方程为\nobreak(\kongwei)

\bindopt {\fontsize{10.5pt}{\qpTJgLeadLiti}\selectfont \makebox[\dimexpr0.25\linewidth-0.25\lxhang\relax][l]{A．\(x=2\)}\makebox[\dimexpr0.25\linewidth-0.25\lxhang\relax][l]{B．\(y=2\)}\makebox[\dimexpr0.25\linewidth-0.25\lxhang\relax][l]{C．\(x=3\)}\makebox[\dimexpr0.25\linewidth-0.25\lxhang\relax][l]{D．\(x=6\)}\par}

\begin{ansblock}[2章-练-课时04-E11]
% ans:2章-练-课时04-E11
\ansitem{11}{B}
\ansnote{详解}{\(M\)、\(N\)纵坐标相等（都为2），直线\(MN\)与\(x\)轴平行，斜率为0，方程为\(y=2\)（同纵坐标的两点连线与\(x\)轴平行）．选：B．}
\end{ansblock}

\tihao{12}\zu{直线方程形式与互化×单选} \tieside{简单(知识点二)}（预习注：本题部分详解使用课时05将学的截距式 \(x/a+y/b=1\)；可先按预习研读，正式学习见课时05）过点\(A(1,4)\)，且横、纵截距的绝对值相等的直线共有\nobreak(\kongwei)

\bindopt {\fontsize{10.5pt}{\qpTJgLeadLiti}\selectfont \makebox[\dimexpr0.25\linewidth-0.25\lxhang\relax][l]{A．1条}\makebox[\dimexpr0.25\linewidth-0.25\lxhang\relax][l]{B．2条}\makebox[\dimexpr0.25\linewidth-0.25\lxhang\relax][l]{C．3条}\makebox[\dimexpr0.25\linewidth-0.25\lxhang\relax][l]{D．4条}\par}

\begin{ansblock}[2章-练-课时04-E12]
% ans:2章-练-课时04-E12
\ansitem{12}{C}
\ansnote{详解}{分类：①过原点（两截距均为0）：\(y=4x\)，1条；\par\noindent ②不过原点，设\(x/a+y/b=1\)，\(|a|=|b|\neq0\)：代入\((1,4)\)得\(1/a+4/b=1\)．\(b=a\)时\(5/a=1\)，\(a=5\)（\(x+y=5\)）；\(b=-a\)时\(1/a-4/a=1\)，\(a=-3\)（\(-x/3+y/3=1\)，即\(y-x=3\)）．共2条．\par\noindent 合计3条．选：C．}
\end{ansblock}

\tihao{13}\zu{直线方程形式与互化×单选} \tieside{简单(知识点二)}（预习注：本题详解同判括注使用课时05将学的截距式 \(x/a+y/b=1\)；可先按预习研读，正式学习见课时05）直线\(5x-2y-10=0\)在\(x\)轴上的截距为\(a\)，在\(y\)轴上的截距为\(b\)，则有\nobreak(\kongwei)
\bindopt {\fontsize{10.5pt}{\qpTJgLeadLiti}\selectfont \makebox[\dimexpr0.5\linewidth-0.5\lxhang-0.25em\relax][l]{A．\(a=2\)，\(b=5\)}\makebox[\dimexpr0.5\linewidth-0.5\lxhang-0.25em\relax][l]{B．\(a=2\)，\(b=-5\)}\par}
\bindopt {\fontsize{10.5pt}{\qpTJgLeadLiti}\selectfont \makebox[\dimexpr0.5\linewidth-0.5\lxhang-0.25em\relax][l]{C．\(a=-2\)，\(b=5\)}\makebox[\dimexpr0.5\linewidth-0.5\lxhang-0.25em\relax][l]{D．\(a=-2\)，\(b=-5\)}\par}

\begin{ansblock}[2章-练-课时04-E13]
% ans:2章-练-课时04-E13
\ansitem{13}{B}
\ansnote{详解}{令\(y=0\)得\(x=2\)；令\(x=0\)得\(y=-5\)．即\(a=2\)，\(b=-5\)（化截距式\(x/2+y/(-5)=1\)同判）．选：B．}
\end{ansblock}

\tihao{14}\zu{直线方程形式与互化×填空} \tieside{中档(知识点二)}（预习注：本题部分详解使用课时05将学的截距式 \(x/a+y/b=1\)；可先按预习研读，正式学习见课时05）过点\(P(4,1)\)作直线\(l\)分别交\(x\)轴、\(y\)轴正半轴于\(A\)，\(B\)两点，\(O\)为坐标原点．当\(|OA|+|OB|\)取最小值时，直线\(l\)的方程为\kongda{x＋2y−6＝0}．

\begin{ansblock}[2章-练-课时04-E14]
% ans:2章-练-课时04-E14
\ansitem{14}{x＋2y−6＝0}
\ansnote{详解}{设\(l:x/a+y/b=1\)（\(a>0\)，\(b>0\)），代入\(P(4,1)\)：\(4/a+1/b=1\)．\par\noindent\(|OA|+|OB|=a+b=(a+b)(4/a+1/b)=5+4b/a+a/b\geqslant5+2\sqrt{(4b/a)\cdot(a/b)}=9\)，当且仅当\(4b/a=a/b\)即\(a=2b\)时取等；联立\(4/a+1/b=1\)代入\(a=2b\)：\(4/(2b)+1/b=1\)\(\Rightarrow3/b=1\)\(\Rightarrow b=3\)，\(a=6\)．\par\noindent 此时\(l\)为\(x/6+y/3=1\)，即\(x+2y-6=0\)．}
\end{ansblock}

% ============================================================
% 思维探索（题15~16＝E15~E16：解答2，前移迁入席居末档照库序）
% ============================================================
\dangbadge{思}{维}{探}{索}

\tihao{15}\zu{直线方程形式与互化×解答} \tieside{简单(知识点三)}分别求下列直线的方程：(1) 经过\(A(0,-1)\)，\(B(-1,1)\)的直线\(l_1\)； (2) 经过\(C(3,-1)\)，\(D(-2,-1)\)的直线\(l_2\)； (3) 经过\(E(1,3)\)，\(F(1,-2)\)的直线\(l_3\)．
\liubai[32mm]{解答书写区}

\begin{ansblock}[2章-练-课时04-E15]
% ans:2章-练-课时04-E15
\ansitem{15}{(1) y＝−2x−1；(2) y＝−1；(3) x＝1}
\ansnote{详解}{(1) 两点横、纵坐标均不等，斜率\(k=(1+1)/(-1-0)=-2\)，点斜式（用\(A\)）：\(y+1=-2(x-0)\)，即\(y=-2x-1\)；\par\noindent (2) \(C\)、\(D\)纵坐标同为\(-1\)，\(l_2\parallel x\)轴，方程为\(y=-1\)（纵坐标相同，与\(x\)轴平行）；\par\noindent (3) \(E\)、\(F\)横坐标同为1，\(l_3\perp x\)轴，方程为\(x=1\)（横坐标相同，为竖直线\(x=1\)）．}
\end{ansblock}

\tihao{16}\zu{直线方程形式与互化×解答} \tieside{中档(知识点二)}（预习注：本题(2)问详解使用课时05将学的截距式 \(x/a+y/b=1\)；可先按预习研读，正式学习见课时05）根据下列条件，求直线的方程：(1) 过点\(A(-3,2)\)，且（在\(y\)轴上的）截距是\(-2\)；(2) 过点\(A(3,0)\)，且在两坐标轴上的截距和为5．
\liubai[24mm]{解答书写区}

\begin{ansblock}[2章-练-课时04-E16]
% ans:2章-练-课时04-E16
\ansitem{16}{(1) y＝−(4/3)x−2；(2) x/3＋y/2＝1（即2x＋3y−6＝0）}
\ansnote{详解}{(1) 截距\(b=-2\)，直线过\((0,-2)\)与\((-3,2)\)，斜率\(k=(2+2)/(-3-0)=-4/3\)，斜截式\(y=-(4/3)x-2\)；\par\noindent (2) 直线过\((3,0)\)，故\(x\)轴截距\(a=3\)（非零），\(y\)轴截距\(b=5-3=2\)（非零），由截距式\(x/3+y/2=1\)，即\(2x+3y-6=0\)．}
\end{ansblock}

% ---- 尾块（\tailfill 置 \end{multicols} 前；无参＝内容尾随框，每件恰 1 处） ----
\tailfill

\end{multicols}
\end{document}
